\documentclass{article}
\usepackage{amsmath,amssymb,hyperref}
\title{The two-level native state hom and its finite-list normalization}
\date{15 November 2026}
\begin{document}\maketitle
\paragraph{Status.} Independently reviewed hom-level construction, extending the finite-depth flattening result. The input is the bicategorical process $K$ and coherent strong pure-state action of a symmetric monoidal bicategory $S$. The first horizontal native construction is assumed established under those hypotheses. This note constructs a two-level \emph{hom} and compares it with the two-state-list dialgebra. It does not assume a second globally assembled state tricategory or its monad laws.

\paragraph{Two-level data with their actual dimension.} At interfaces $A,B$, an object is an outer state $X$ and an inner machine $(Y,f)$ whose body has type
\[
 f:L_Y(L_XA)\longrightarrow L_Y(L_XB).
\]
Action associativity identifies these endpoints with $L_{Y\otimes X}A,B$ by chosen equivalences; the body is conjugated by those equivalences when changing presentations. An outer native arrow consists of $h:X\to X'$ and an inner native arrow between the composites with the unit-state pure boundary machines induced by $h$. The two inner state objects in those composites are $Y\otimes I$ and $I\otimes Y'$ (or the reverse unitor convention if composition notation is reversed). It therefore includes a state map
\[
 \widehat k:Y\otimes I\longrightarrow I\otimes Y'
\]
and a $K$-two-cell. Use the specified unitors and their adjoint inverses to transport $\widehat k$ to $k:Y\to Y'$. Its process cell, transported through the action comparisons, is exactly
\[
 \alpha:L_B(k\otimes h)f\Longrightarrow gL_A(k\otimes h).
\]
This statement uses strong pure centrality to identify the two boundary pastes; neither $f$ nor $g$ is moved past another process.

An outer three-cell is a pair: $\sigma:h\Rightarrow h'$ and an inner native two-cell $\widehat\tau$ between the appropriate inner boundary-pasted state maps. Transporting through unitors and the pure-action pseudonaturality gives $\tau:k\Rightarrow k'$. The inner native equation is the $K$-two-cell equality
\[
 \bigl(gL_A(\tau\otimes\sigma)\bigr)\alpha
 =\alpha'\bigl(L_B(\tau\otimes\sigma)f\bigr),
\]
including the tensor/action compositors. Thus the extra three-cell witness has not been discarded: it is exactly the second state two-cell $\tau$ with its native equation. Forgetting only $\sigma$ would lose data.

\paragraph{Normalized hom construction.} Define $N_2(A,B)$ directly by the tuples $(X,Y,f)$, arrows $(h,k,\alpha)$, and two-cells $(\sigma,\tau)$ with the displayed equation. Compose the state tuples in $S\times S$, and compose their native cells with the pseudofunctor comparisons of $m_2(X,Y)=Y\otimes X$. This is precisely
\[
 D(L_A m_2,L_B m_2)=D_2(A,B).
\]
Consequently identities, composition, hom associators/unitors, and two-cell interchange exist and satisfy their equations by the established dialgebra construction. In particular this is not a construction by simply forgetting the inner three-cell.

\paragraph{Comparison with the unnormalized two-level hom.} One can equivalently keep $\widehat k,\widehat\tau$ and the unit-state boundary composites. Define composition by their actual first-level pasting. Normalization conjugates each of the two intermediate unit states by the chosen unitors. This is an equivalence on the corresponding state hom categories; its inverse inserts those unitors. On process cells use the same conjugation and action comparisons, so the native-square condition is preserved in both directions. Units/counits are induced by the chosen unitor adjoint equivalences. For two composable arrows the comparison between normalizing their paste and composing their normalizations is the pseudonaturality/composition constraint of those unitors and the pure action. Its admissibility is the pasting of the two given native squares; naturality at $(\sigma,\tau)$ is the corresponding native equation. The comparison pentagon and units are the unitor/action coherence equations. This gives a biequivalence of the two hom presentations, rather than literal equality of raw boundary maps.

The scope of this assertion is fixed interfaces and two state layers. It does not prove that a freely chosen tricategorical lax-comma construction has the same global composition, or that arbitrary process-three-cell witnesses reduce to state pairs. The reduction applies because the process tricategory at the second level is specifically the first native state construction on a bicategory $K$.

\paragraph{Actual local flattening.} Composing the normalization with the reviewed $F_2:D_2(A,B)\to D_1(A,B)$ gives a local flattening pseudofunctor. It sends $(X,(Y,f))$ to $(Y\otimes X,\bar f)$, the pair of state maps to $k\otimes h$, and the native three-cell witness pair to $\tau\otimes\sigma$. Its comparison cells are explicit tensor/action comparisons. Tensor need not be faithful and is not used to detect an unproved source equation. The image equation follows by applying the actual tensor/action pseudofunctors to the already satisfied native equation.

\paragraph{Remaining global proof.} Horizontal compatibility requires the middle-coordinate permutation between $(Y\otimes X)\otimes(V\otimes U)$ and $(Y\otimes V)\otimes(X\otimes U)$, together with the same chronological bodies. Its naturality must be checked with both native witness coordinates, and its comparison modifications must retain their underlying state two-cells. Depth-three and depth-four multiplication associativity then require an independently defined global second state constructor and its structural transformation doctrine. The present local biequivalence removes the first typing/dimension obstruction but does not assert those additional constructions.
\end{document}
